%% =====================================================================
%% Paper U — Discriminant topology and the critical points of a
%% spectral rigidity functional on conformal deformations of the
%% pillowcase
%% Author: Dhiren J. Master
%% Build: v4, 11 July 2026. Standalone document; BF register.
%% v2: landing paragraph after Theorem 1.1 + echo in proof of Thm A.
%% v3: GD established at reduction level (Proposition gd:verdict,
%% GD_certificate.lean); GD retired from all displayed hypothesis
%% sets; section 7.2 consumption sentence; abstract strengthened.
%% v4: O4 readout executed (O4_readout_certificate.lean) — counting
%% object identified as the fibre Euler characteristic; flattening
%% exclusivity lemma; section 7.2 upgraded from test to result;
%% section 9 amended.
%% Rulings consumed: R1 both routes (Route M spine); R2 cite Paper 3
%% only, Paper 5 items reproduced; R3 JGP positioning; R4 doubling
%% reproduced in-document; R5 GD discharged in-session, O4 readout
%% displayed; R6
%% title per PI.
%% =====================================================================
\documentclass[11pt]{article}

\usepackage[a4paper, margin=2.9cm]{geometry}
\usepackage{amsmath, amssymb, amsthm}
\usepackage{array}
\usepackage[british]{babel}
\usepackage[hidelinks]{hyperref}

\input{mef_macros.tex}

\numberwithin{equation}{section}

\theoremstyle{plain}
\newtheorem{theorem}{Theorem}[section]
\newtheorem{proposition}[theorem]{Proposition}
\newtheorem{lemma}[theorem]{Lemma}
\newtheorem{corollary}[theorem]{Corollary}
\newtheorem{hypothesis}[theorem]{Hypothesis}

\theoremstyle{definition}
\newtheorem{definition}[theorem]{Definition}

\theoremstyle{remark}
\newtheorem{remark}[theorem]{Remark}

\newcommand{\Zn}[1][n]{\mathbb{Z}_{#1}}
\newcommand{\stildec}{\tilde{\sigma}}
\newcommand{\Xad}{X_{\mathrm{ad}}}
\newcommand{\Xbad}{\bar{X}_{\mathrm{ad}}}

\title{Discriminant topology and the critical points of a spectral
rigidity functional on conformal deformations of the pillowcase}
\author{Dhiren J. Master\\[2pt]
\normalsize Independent Researcher\\[1pt]
\normalsize \texttt{dhiren@masterequation.org}}
\date{}

\begin{document}
\maketitle

\begin{abstract}
\noindent
The discriminant of a family of complex Hermitian operators is,
after von Neumann and Wigner, a locus of codimension three, and the
complement of such a locus acquires its topology from the
two-spheres that link it. We study the consequences of this
principle for a spectral rigidity functional on a space of conformal
deformations of the pillowcase orbifold $T^{2}/\Ztwo$. Within a
smoothed spectral window we enumerate the discriminant strata,
establish a parity selection rule, a channel-gating non-degeneracy
and a Gram non-degeneracy of the adjacent exchanges, and prove that
the admissible deformation space
carries homology of total rank twelve at generator level, under a
hypothesis set displayed in full. The Morse inequalities for the
functional are then equalities, so that critical points are counted
by Betti numbers; in particular the functional admits exactly one
stable critical point. An independent route derives the same
uniqueness from two local hypotheses alone. The topology of the
discriminant complement thereby reduces the rigidity question to a
displayed set of structural hypotheses together with finite
non-degeneracy checks, the latter established here with machine
verification.
\end{abstract}

\section{Introduction}\label{sec:intro}

The purpose of this paper is to determine the critical structure of
a spectral rigidity functional on a space of conformal deformations
of the pillowcase orbifold, by computing the topology of the space
on which the functional lives. The mechanism is classical. Von
Neumann and Wigner \cite{vonNeumannWigner1929} observed that a
coincidence of two eigenvalues of a complex Hermitian operator
imposes three real conditions, so that within a generic family the
discriminant --- the locus of degenerate spectrum --- has
codimension three; Berry \cite{Berry1984} and Simon
\cite{Simon1983} identified the two-spheres linking such a locus as
carriers of a unit Chern charge. A functional defined through the
spectrum is smooth away from the discriminant and, in general, not
through it; the natural domain of the variational problem is the
discriminant complement, and whatever homology that complement
carries is imposed on the critical structure of every admissible
functional through the Morse inequalities. The present paper
carries this programme out in full for one specific geometry, from
the enumeration of the discriminant strata to the count of critical
points.

The geometry is the following. The fibre is the pillowcase
$F = \Ttwo/\Ztwo$ at the modulus $\tau = i\tautwo$,
$\tautwo = 93/50$, decorated with the $\Spinc$ structure of
determinant line bundle $L = \Oline{3}$; the base is
$B = \CPtwo \times \Stwo$; and the total space is the warped
product
\begin{equation}\label{eq:K8}
\Keight \;=\; \bigl[(\CPtwo \times \Stwo) \times_{w}
(\Ttwo/\Ztwo)\bigr]^{\Spinc},
\qquad \chi(\Keight) = 12 .
\end{equation}
The functional, written $\omega$ throughout, is the Dyson--Mehta
spectral rigidity statistic \cite{DysonMehta1963} of the $\Spinc$
Dirac spectrum, computed over a smoothed spectral window and
regarded as a functional of a conformal deformation of the fibre
metric, modulated over the base and constrained to fixed fibrewise
area. The variational theory of this functional on the fibre alone
--- criticality of the constant configurations, the
density-balance form of the Euler--Lagrange equation, and the
residual symmetry of the problem --- was established in
\cite{MEF3}, and those results are consumed here at their stated
grades; everything concerning the base-modulated enlargement, the
discriminant of the assembled operator, and the topology of its
complement is proved in the present document.

Our main results are the following. The hypotheses they consume
are displayed in full where they are first used, each stated once
with what it suppresses named; the introduction records only their
labels.

\begin{theorem}[Generator-level topology of the admissible
region]\label{thm:A}
Assume the smoothed window weight of
Definition~\textup{\ref{rb2:window}}, the retention--detection
statement OW \textup{(Hypothesis~\ref{rb2:ow})}, the transgression
compatibility TR \textup{(Hypothesis~\ref{rb2:tr})}, and the
higher-strata dependence HS \textup{(Hypothesis~\ref{hyp:hs})}; the
channel-gating non-degeneracy NDC and the Gram non-degeneracy GD
hold at reduction level on the verified list
\textup{(Propositions~\ref{ndc:verdict} and~\ref{gd:verdict})}.
Then there is a finite
$\varepsilon^{*}$ such that for every
$\varepsilon \geq \varepsilon^{*}$ the rational cohomology of the
admissible adiabatic region is the truncated polynomial ring
\[
H^{*}\bigl(\Xad(\varepsilon); \Q\bigr) \;\cong\;
\Q[x, y, z]\,/\,(x^{3},\, y^{2},\, z^{2}),
\qquad |x| = |y| = |z| = 2,
\]
with Betti sequence $(1, 0, 3, 0, 4, 0, 3, 0, 1)$, of total rank
twelve; and the inclusions of the truncations into the full
base-modulated deformation space are rational homology
isomorphisms, so the same sequence is carried by the unrestricted
space.
\end{theorem}

The ring displayed is the rational cohomology ring of $\Keight$
itself, generator for generator and relation for relation
(Proposition~\ref{prop:betti}): given Theorem~\ref{thm:A}, the
discriminant complement of the deformation problem is a
rational-cohomology copy of the manifold being deformed, and the
Morse theory of $\omega$ is, at the level of homology, the Morse
theory of a function on $\Keight$. Every count in this paper is a
consequence of that identification.

\begin{theorem}[The critical count]\label{thm:B}
Assume, in addition to the hypotheses of
Theorem~\textup{\ref{thm:A}}, the Morse-theoretic admissibility PS
\textup{(Hypothesis~\ref{o2:ps})} and the upper bound UB
\textup{(Hypothesis~\ref{hyp:ub})}. Then the critical points of
$\omega$ are counted exactly by the Betti numbers,
$c_{i} = b_{i}$ for every index $i$; in particular
\[
c_{0} \;=\; 1 :
\]
the functional admits exactly one stable critical point, and the
remaining critical structure consists of $3 + 4 + 3 + 1$ saddles
of even index $2, 4, 6, 8$.
\end{theorem}

\begin{theorem}[Uniqueness by the local route]\label{thm:C}
Assume instead the two local hypotheses B
\textup{(Hypothesis~\ref{o1:break})} and U
\textup{(Hypothesis~\ref{o1:symuniq})}. Then $c_{0} = 1$, with no
input from Theorem~\textup{\ref{thm:A}} and no claim at higher
index.
\end{theorem}

\begin{corollary}\label{cor:main}
Under either displayed hypothesis set, $\omega$ admits exactly one
critical point of index zero.
\end{corollary}

The shape of the argument is as follows. The ambient deformation
spaces are shown to be contractible by an explicit retraction built
from the heat semigroup on the base composed with a fibrewise area
renormalisation, so that every homology class of the admissible
region is borne by the discriminant (Section~\ref{sec:ambient}).
The discriminant of the assembled operator is then enumerated by
the base content of the coinciding pair; a parity selection rule
removes the coincidences of opposite fibre parity, and a gating
computation identifies which of the remaining strata carry the
generic codimension-three structure with charged linking spheres
(Section~\ref{sec:discriminant}). The smoothed window converts the
window edges from codimension-one walls into boundary structure on
the in-window strata, and an edge-capping argument turns that
boundary structure into linear relations among the linking classes;
the relations collapse each base-content sector to a single
generator, and a transgression pairing against the base
intersection forms fixes the ring relations, the generator with
surviving square being necessarily the one over the projective
factor (Section~\ref{sec:collapse}). The same transversal exchanges
that collapse the sectors exhaust the defect of the retraction at
finite deformation energy, which is the stabilisation giving
Theorem~\ref{thm:A} (Section~\ref{sec:stab}). The count then
follows the classical route: the second variation at the constant
stratum decomposes into six blocks labelled by the harmonic
bidegrees of the base, an anisotropic scaling argument separates
the six dressed weights pairwise, a group-replacement family
$\Ttwo/\Zn$ separates the three candidate values of the fibre
counting factor and identifies the counting object as the fibre
Euler characteristic, and the Morse inequalities, having no odd
Betti numbers to accommodate, collapse to equalities under the
displayed count hypothesis (Section~\ref{sec:count}). The local route of
Theorem~\ref{thm:C} is independent of all of this and occupies
Section~\ref{sec:retreat}.

A number of finite verifications --- monomial counts, dimension
arithmetic, sign and parity bookkeeping, the reference level list
and its crossing pairs --- are machine-checked in Lean~4; the
certificates are cited where they arise and are available from the
author. The checks are elementary in every case and the reader can
reproduce each by hand; the certificates exist so that no
bookkeeping step of the argument rests on care alone.

The paper is organised as follows. Section~\ref{sec:setting} fixes
the geometry, the spectrum, the functional and the configuration
spaces, and records the results of \cite{MEF3} that are consumed.
Section~\ref{sec:ambient} proves the contractibility of the ambient
spaces and the colimit identity. Section~\ref{sec:discriminant}
enumerates the discriminant, proves the parity selection rule and
the gating lemma, and establishes the non-degeneracy condition NDC
at reduction level. Section~\ref{sec:collapse} proves the edge
capping and sector separation lemmas, establishes the Gram
non-degeneracy GD on the verified list, displays the hypotheses
OW and TR, and derives the collapse and the ring structure.
Section~\ref{sec:stab} proves stabilisation and assembles
Theorem~\ref{thm:A}. Section~\ref{sec:count} treats the six
base-sector replicas, the fibre counting object, and the Morse
consumption, proving Theorem~\ref{thm:B}.
Section~\ref{sec:retreat} proves Theorem~\ref{thm:C}.
Section~\ref{sec:outlook} records what remains open and in what
form.

\section{The geometry, the spectrum and the
functional}\label{sec:setting}

This section fixes notation and assembles the objects. Throughout,
$\zeta = \chi(\Keight)\,\pi = 12\pi$ denotes the fixed conformal
normalisation; rational coefficients are understood for all
homology and cohomology.

\subsection{The fibre spectrum and the total space}
\label{sec:fibrespec}

On the torus $\Ttwo$ of modulus $\tau = i\tautwo$ the $\Spinc$
Dirac operator with determinant line bundle $L = \Oline{3}$ has
spectrum labelled by the shifted momentum lattice: the labels are
$(m,\, n - \tfrac{3}{2})$ with $(m, n) \in \Z^{2}$, the
half-integer shift $-c_{1}(L)/2 = -3/2$ contributed by the
$\Spinc$ connection, and
\begin{equation}\label{eq:spectrum}
\lambda^{2}(m, n) \;\propto\;
\frac{m^{2}}{\tautwo^{2}} + \Bigl(n - \frac{3}{2}\Bigr)^{2}.
\end{equation}
At $\tautwo = 93/50$ the squared levels are proportional to the
integer list
\begin{equation}\label{eq:levellist}
V(m, n) \;=\; 10^{4}\,m^{2} + 8649\,(2n - 3)^{2}
\end{equation}
in units of $(2\pi)^{2}/34596$. Because $c_{1}(L) = 3$ is odd, the
involution $\sigma$ acts freely on the shifted lattice: no label is
$\sigma$-fixed, every level of the pillowcase quotient is a genuine
two-element $\sigma$-orbit, and the $\Spinc$ lift $\stildec$ of
$\sigma$ projects the torus problem to the pillowcase problem
without fixed labels \cite{MEF3}. The pillowcase fibre is
$F = \Ttwo/\Ztwo$ with this decoration, the base is
$B = \CPtwo \times \Stwo$, and the total space is the warped
product \eqref{eq:K8}, the warp acting on the fibre metric over the
base. We record the homology of the total space at once; the
computation is elementary and is used only through its output.

\begin{proposition}[The Betti sequence]\label{prop:betti}
The rational Betti numbers of $\Keight$ are
\begin{equation}\label{eq:betti}
b(\Keight) \;=\; (1,\, 0,\, 3,\, 0,\, 4,\, 0,\, 3,\, 0,\, 1),
\end{equation}
of total rank $12 = \chi(\Keight)$.
\end{proposition}

\begin{proof}
The rational cohomology of the pillowcase is that of the sphere:
the involution acts as $-1$ on $H^{1}(\Ttwo; \Q) = \Q^{2}$, so the
invariant part vanishes, giving Betti numbers $(1, 0, 1)$. The
factors $\CPtwo$ and $\Stwo$ contribute $(1, 0, 1, 0, 1)$ and
$(1, 0, 1)$. All factors have vanishing odd cohomology, so the
rational Leray spectral sequence of the warped product degenerates
and the K\"unneth formula applies; the Poincar\'e polynomial is
\[
(1 + t^{2} + t^{4})(1 + t^{2})(1 + t^{2})
\;=\; 1 + 3t^{2} + 4t^{4} + 3t^{6} + t^{8},
\]
which is \eqref{eq:betti}. With all odd $b_{i}$ zero the
alternating and plain sums coincide at $12$; Poincar\'e symmetry is
manifest. The monomial count realising the sequence is
machine-verified by two independent routes (certificates
\texttt{RB1\_certificate.lean}, \texttt{RB2\_certificate.lean}).
\end{proof}

\subsection{The rigidity functional and the smoothed window}
\label{sec:functional}

Let $\{\lambda_{k}\}$ be the spectrum of the $\Spinc$ Dirac
operator of a configuration, listed with multiplicity, and let
$x_{k}$ denote the unfolded positions: the images of the
$\lambda_{k}$ under the mean counting function, so that the
unfolded sequence has unit mean spacing. The Dyson--Mehta statistic
$\Delta_{3}(L)$ of a window of unfolded length $L$ is the mean
square deviation of the unfolded staircase from its optimal linear
fit over the window \cite{DysonMehta1963}; it is the classical
measure of spectral rigidity, minimal for crystalline spectra and
maximal for Poissonian ones. The functional studied in this paper
is the smoothed windowed form of this statistic.

\begin{definition}[The smoothed window and the
functional]\label{rb2:window}
Fix $\hat{w} \in C^{\infty}_{c}(\R)$ with
$\operatorname{supp}\hat{w} = [a, a+L]$ in the unfolded variable,
$0 < \hat{w} \leq 1$ on the interior, and all derivatives vanishing
at the endpoints. The functional $\omega$ is the rigidity statistic
computed with weights $w_{k}\,\hat{w}(x_{k})$, where
$w_{k} = \tfrac{2}{L}\hat{D}(x_{k})\,\bar\rho(\lambda_{k})$ are the
density-balance weights, $\hat{D}$ the deviation of the unfolded
staircase from its optimal linear fit and $\bar\rho$ the mean
density, as in the criticality equation of \cite{MEF3}.
\end{definition}

The fibre-only variational theory of $\omega$ is that of
\cite{MEF3} and we consume three of its results without change.
First, on the fixed-area slice of $\sigma$-even conformal
deformations of the fibre, the constant configurations are critical
points of $\omega$ for every window, degenerate blocks included,
and the criticality is symmetry-protected: it holds for every
choice of rigidity weights, so nothing below depends on the
one-hump structure of $\Delta_{3}$. Second, on the simple stratum
the criticality condition takes the density-balance form
\begin{equation}\label{eq:el}
\sum_{k \in W} w_{k}\,\lambda_{k}\,
\bigl|\chi_{k}^{\Phi}(x)\bigr|^{2}
\;=\; \mu\, e^{2\zeta\Phi(x)}
\qquad \text{for every } x \in F,
\end{equation}
with $\chi_{k}^{\Phi}$ the eigenspinors of the deformed operator
and $\mu$ the area-constraint multiplier. Third, of the half-period
translations of the torus only the unshifted one survives the
$\Spinc$ projection --- the two shifted half-periods conjugate the
lift $\stildec$ with sign $-1$, precisely because $c_{1}(L) = 3$ is
odd --- so the problem carries a residual $\Ztwo$ acting by
$\omega$-preserving isometries, and by Palais' principle of
symmetric criticality \cite{Palais1979} every critical point of the
restriction to the symmetric slice is a critical point of $\omega$.
We write $X^{\Ztwo}$ for the symmetric slice throughout. The
compatibility of the smoothed weight with these consumed results is
immediate: the symmetry-protected criticality holds for every
weight system, the density-balance equation holds with the dressed
weights $w_{k}\hat{w}(x_{k})$ in place of $w_{k}$, and the
perturbative inputs (Kato regularity on the simple stratum) are
untouched.

\subsection{The configuration spaces}\label{sec:spaces}

\begin{definition}[The base-modulated configuration
space]\label{rb:space}
The base-modulated configuration space $\bar{X}_{B}$ is the Banach
manifold of $C^{k}$ profiles $\Phi : B \times F \to \R$ that are
$\sigma$-even in the fibre argument for every base point and
satisfy the fibrewise area constraint
$\int_{F} e^{2\zeta\Phi(b,\,\cdot)}\, dv_{0} = A_{0}$ for every
$b \in B$. The fibre-only space $\bar{X}_{f}$ of the window theory
is the slice of base-constant profiles, and the results of
\cite{MEF3} are consumed unchanged on that slice. Since the
conformal normalisation and the warp exponent are the same
quantised constant $\zeta$, a base-modulated profile is a warp-type
deformation of $\Keight$ at the flux-quantised exponent: the
enlargement adds no channel the geometry does not already carry.
\end{definition}

\begin{definition}[The adiabatic truncation]\label{rb:trunc}
For $\Phi \in \bar{X}_{B}$ write
\[
E_{B}[\Phi] \;=\; \int_{B \times F}
\bigl|\nabla_{\!B}\,\Phi\bigr|^{2}_{g_{B}}\; dv ,
\]
the base Dirichlet energy, and let $\gamma = 2\pi/\tautwo$ be the
fibre gap of the reference spectrum. The adiabatic space at
parameter $\varepsilon > 0$ is
$\Xbad(\varepsilon) = \{\Phi \in \bar{X}_{B} :
E_{B}[\Phi] \leq \varepsilon^{2}\gamma^{2}\}$. The Dirichlet form
is chosen over any pointwise bound because it is monotone under the
base heat flow, so the truncation is compatible with the retraction
it is about to meet.
\end{definition}

\begin{definition}[The discriminant and the admissible
regions]\label{rb:adm}
Let $\Sigma_{B} \subset \bar{X}_{B}$ be the discriminant of the
assembled problem: the zero-mode, window-edge and in-window
crossing loci of the full $\Keight$ operator over the
base-modulated profile. For a crossing component $\Sigma_{c}$ write
$x_{c}(\Phi)$ for the unfolded position of the coinciding pair; the
in-window part is
$\Sigma_{c}^{w} = \Sigma_{c} \cap \{x_{c} \in (a, a+L)\}$. The
admissible regions are taken with the in-window parts only:
\[
\Xad(\varepsilon) \;=\; \Xbad(\varepsilon) \setminus
\Bigl(\,\bigcup_{c}\Sigma_{c}^{w} \;\cup\; \Sigma_{0}\Bigr),
\qquad
X_{B} \;=\; \bar{X}_{B} \setminus
\Bigl(\,\bigcup_{c}\Sigma_{c}^{w} \;\cup\; \Sigma_{0}\Bigr),
\]
with $\Sigma_{0}$ the zero-mode locus. That this excision, and not
the larger one including the window-edge loci, is the correct
domain of the variational problem is Lemma~\ref{rb2:adm} below. We
write $X$ for $X_{B}$ where no confusion arises.
\end{definition}

\section{The ambient spaces carry no topology}\label{sec:ambient}

The purpose of this section is to prove that $\bar{X}_{B}$ and
every truncation $\Xbad(\varepsilon)$ are contractible, and to
record the colimit identity relating the admissible truncations to
the full admissible space. The consequence used throughout is that
whatever homology the admissible region carries comes from the
discriminant alone.

\begin{definition}[The retraction map]\label{rb:retraction}
Let $H_{t} = e^{t\Delta_{B}}$ act on the base argument of a
profile, the fibre argument carried as a parameter, and let $N$ be
the fibrewise area renormalisation
\[
N(\Phi)(b, \cdot) \;=\; \Phi(b, \cdot) - \tfrac{1}{2\zeta}
\log\!\Bigl(\tfrac{1}{A_{0}}\int_{F}
e^{2\zeta\Phi(b,\,\cdot)}\,dv_{0}\Bigr).
\]
The retraction is the composite $R_{t} = N \circ H_{t}$.
\end{definition}

\begin{lemma}[Heat flow on the base argument]\label{rb:heat}
$H_{t}$ preserves fibrewise $\sigma$-evenness, fixes base-constant
profiles, decreases $E_{B}$ monotonically, and converges as
$t \to \infty$ to the base mean
$\bar\Phi(x) = \textstyle\frac{1}{\Vol(B)}\int_{B}\Phi(b, x)\,dv_{B}(b)$, exponentially at
the rate of the first non-zero eigenvalue of $\Delta_{B}$.
\end{lemma}

\begin{proof}
The flow acts on the base argument and $\sigma$ on the fibre
argument, so the two commute, and evenness, a linear condition, is
preserved by the linear flow. Constants in the base are fixed
points of $H_{t}$. Monotonicity is the standard energy identity
$\tfrac{d}{dt}E_{B}[H_{t}\Phi] = -2\!\int|\Delta_{B}H_{t}\Phi|^{2}
\leq 0$, and convergence to the mean at rate $\lambda_{1}(B)$ is
the spectral decomposition of the heat semigroup on the closed
base. A discrete-step instance of the contraction, with explicit
rational eigendata and the exact energy-drop identity, is
machine-verified (certificate
\texttt{RB\_retraction\_certificate.lean}).
\end{proof}

\begin{lemma}[Jensen one-sidedness of the fibre
areas]\label{rb:jensen}
Along the flow the fibre areas can only shrink:
$A_{t}(b) = \int_{F} e^{2\zeta(H_{t}\Phi)(b,\,\cdot)}\,dv_{0} \leq
A_{0}$ for every $b$ and $t$, since $H_{t}\Phi(b, x)$ is a convex
average of $\Phi(\cdot, x)$ over base points and the exponential is
convex. Consequently the renormalisation shift
$c_{t}(b) = -\tfrac{1}{2\zeta}\log(A_{t}(b)/A_{0})$ is
non-negative, continuous, and constant on each fibre, so $N$ is
well-defined on the flow image and preserves fibrewise
$\sigma$-evenness.
\end{lemma}

\begin{lemma}[Quasi-invariance of the truncation]\label{rb:quasi}
The composite satisfies
$R_{t} : \Xbad(\varepsilon) \to \Xbad(C\varepsilon)$ for all
$t \geq 0$, with a constant $C$ controlled as follows. The base
gradient of the shift obeys the chain rule
$\nabla_{\!B}\, c_{t}(b) = -\!\int_{F}
\nabla_{\!B}(H_{t}\Phi)(b, x)\, d\mu_{b}(x)$ with
$d\mu_{b} = e^{2\zeta H_{t}\Phi}\,dv_{0}/A_{t}(b)$ a probability
measure on the fibre, so by Cauchy--Schwarz
$|\nabla_{\!B}c_{t}|^{2}(b) \leq \int_{F}
|\nabla_{\!B}H_{t}\Phi|^{2}\,d\mu_{b}$; on the fixed-area slice the
deformed and reference fibre measures are uniformly comparable for
profiles of bounded fibre oscillation, and combining with the
monotonicity of Lemma~\ref{rb:heat} gives
$E_{B}[R_{t}\Phi] \leq C^{2} E_{B}[\Phi]$ with $C$ explicit in the
comparability constant. The comparability constant depends on a
fibre-oscillation bound on the slice, which we state rather than
suppress; in the directed system of truncations a bounded loss of
constant is harmless, the maps below being consumed only up to
enlargement of $\varepsilon$.
\end{lemma}

\begin{proposition}[The ambient spaces are
contractible]\label{rb:ambient}
$\bar{X}_{B}$ and every $\Xbad(\varepsilon)$ are contractible: the
flow-with-renormalisation $R_{t}$ deformation retracts each onto
the fibre-only slice through profiles of non-increasing base
energy, and the affine contraction of the fibre-only slice
completes the retraction to a point. Whatever topology the
admissible region carries therefore comes from the discriminant
alone.
\end{proposition}

\begin{lemma}[Colimit identity]\label{rb:colimit}
$X_{B} = \bigcup_{\varepsilon > 0} \Xad(\varepsilon)$ is an
increasing union of open subsets, so singular homology satisfies
$H_{*}(X_{B}) = \varinjlim_{\varepsilon}
H_{*}(\Xad(\varepsilon))$ \cite{Hatcher2002}: every class of the
full space is realised at some finite $\varepsilon$, and no class
is invisible to all truncations.
\end{lemma}

The statement that the homology of the truncations stabilises at a
finite parameter --- so that the truncated and unrestricted
formulations of every result below coincide --- is
Theorem~\ref{rb2:stab}; it is proved in Section~\ref{sec:stab} as a
consequence of the collapse machinery, not assumed.

\section{The assembled discriminant}\label{sec:discriminant}

This section enumerates the discriminant of the assembled operator,
removes the uncharged strata by a parity selection rule, identifies
the charged strata by a gating computation, and establishes the
non-degeneracy of the gating at the lowest levels of both towers.
The output is the input of the collapse machinery of
Section~\ref{sec:collapse}.

\begin{definition}[Strata of the assembled
discriminant]\label{rb1:strata}
At the adiabatic point the spectrum of the assembled operator
carries the block form
$\Lambda^{2} = \lambda_{b}^{2} + \lambda_{f}^{2}(\Phi)$, levels
labelled by a base state and a fibre label. The assembled
discriminant $\Sigma_{B}$ decomposes by the base content of the
coinciding pair: the fibre stratum (equal base data, fibre
reflection partners recombining), the tower strata (one harmonic
and one excited base state, indexed by the base tower and level),
and the residual strata (two excited base states, and transverse
intersections of the foregoing).
\end{definition}

\begin{lemma}[Out-of-window coincidences are
admissible]\label{rb2:adm}
The functional of Definition~\textup{\ref{rb2:window}} is
$C^{\infty}$ at every profile whose only spectral coincidences lie
outside the open window. Consequently the excision of
Definition~\textup{\ref{rb:adm}} is the correct one: no window-edge
locus is removed, and the codimension-one walls of the sharp-window
theory do not arise.
\end{lemma}

\begin{proof}
Away from coincidences the eigenvalue branches are analytic in
$\Phi$ on the simple stratum \cite{Kato} and the composite is
smooth. At a coincidence of unfolded position outside $(a, a+L)$
the branches are merely Lipschitz, but every occurrence of the
non-smooth pair in the statistic is multiplied by $\hat{w}$
evaluated where it vanishes to infinite order together with all
derivatives; the chain rule therefore produces smooth limits of
every order from the window factor, and the contribution of the
pair to each derivative of the functional extends continuously by
zero. The remaining terms involve only spectrally isolated levels
and are smooth as before.
\end{proof}

\begin{lemma}[Parity selection]\label{rb1:parity}
The involution acts on fibre labels alone, the base being
$\sigma$-inert; hence two levels of opposite fibre $\sigma$-parity
belong to orthogonal sectors for every admissible profile, their
coincidences carry no coupling and no eigenbundle charge, and under
the smoothed window weight the functional is smooth through them.
Such coincidences are not discriminant points. The parity
bookkeeping is machine-verified (certificate
\texttt{RB1\_certificate.lean}).
\end{lemma}

\begin{lemma}[Warp gating of the tower strata]\label{rb1:gate}
At the strictly adiabatic point with the warp switched off, the
off-diagonal element between levels of distinct base states
factorises through the base inner product and vanishes identically;
the corresponding coincidences are codimension-one, uncharged, and
carry no linking topology. With the warp present, the excited-block
coupling dresses the fibre deformation with a base-dependent
factor, the off-diagonal element acquires the form of a base
integral against the modulation profile, and wherever this integral
is non-zero the coincidence is a generic complex Hermitian
crossing: codimension three, with linking two-spheres of unit Chern
charge \cite{vonNeumannWigner1929, Berry1984, Simon1983}. The
ground sector, being exactly warp-neutral by the conformal
covariance of the Dirac operator
\cite{Hitchin1974, Hijazi, BourguignonGauduchon1992}, is untouched:
the fibre stratum retains its six-fold simultaneous dressing and
contributes a single class per reflection pair.
\end{lemma}

The non-vanishing of the gating integrals for specific tower levels
is a non-degeneracy condition, held apart from the lemma and now
established at reduction level. The mechanism deserves one sentence
of orientation: because the conformal normalisation and the warp
exponent are the same quantised constant
(Definition~\ref{rb:space}), the gating dressing is supplied by the
base-modulated deformation itself, so the condition is a statement
about deformation directions and not about any background profile.

\begin{lemma}[Factorisation of the gating element]\label{ndc:factor}
Let $(h \otimes f_{k})$ and $(e \otimes f_{k'})$ be assembled
levels of equal fibre $\sigma$-parity, $h$ the constant base state
and $e$ a normalised base tower harmonic, and let
$\delta\Phi(b, x) = Y(b)\,\varphi(x)$ be an admissible product
deformation, $Y$ a base harmonic and $\varphi$ fibrewise
$\sigma$-even. The first-order off-diagonal element of the
assembled operator between the two levels factorises,
\begin{equation}\label{eq:gating}
\bigl\langle e \otimes f_{k'}\,\bigl|\,\delta D\,\bigr|\,
h \otimes f_{k}\bigr\rangle
\;=\;
\langle e \,|\, Y \,|\, h\rangle_{B}\;\cdot\;
\mathcal{F}_{k'k}[\varphi],
\end{equation}
with $\mathcal{F}_{k'k}[\varphi]$ the fibre matrix element of the
conformal perturbation between the shifted-lattice eigenspinors. If
$\varphi$ has no constant Fourier component, the deformation
preserves the fibrewise area constraint to first order.
\end{lemma}

\begin{proof}
The perturbation induced by a base-modulated conformal deformation
acts, to first order in the warped reduction, as multiplication by
$Y$ on the base argument tensored with the fibre conformal
perturbation operator determined by $\varphi$ through the conformal
covariance of the Dirac operator
\cite{Hitchin1974, Hijazi, BourguignonGauduchon1992}; the base and
fibre integrations separate. The perturbative input is standard
first-order theory on the simple stratum \cite{Kato}, evaluated in
the degenerate two-level block at the coincidence. First-order area
preservation for mean-free $\varphi$ is the linearisation of the
fixed-area constraint.
\end{proof}

\begin{lemma}[The base factor is exact and positive]\label{ndc:base}
Choosing $Y = e$, the base factor is
$\langle e \,|\, e \,|\, h\rangle_{B}
= \Vol(B)^{-1/2}\,\|e\|_{L^{2}}^{2} > 0$, for every tower and every
level. The base side of the non-degeneracy condition therefore
holds identically, and the entire content of the condition is
carried by the fibre factor.
\end{lemma}

\begin{proof}
$h = \Vol(B)^{-1/2}$ is constant, so the element is
$\Vol(B)^{-1/2}\!\int_{B} e^{2}\,dv_{B}$, which is the displayed
norm. Positivity is immediate.
\end{proof}

\begin{lemma}[The fibre matrix element on the shifted
lattice]\label{ndc:fibre}
On the flat pillowcase at $\tau = i\tautwo$ the eigenspinors carry
momenta $\kappa(m, n) \propto (m/\tautwo,\; n - \tfrac{3}{2})$ on
the shifted lattice, and for the $\sigma$-projected states the
fibre element is the finite channel sum
\[
\mathcal{F}_{k'k}[\varphi] \;=\;
c(\lambda_{k}, \lambda_{k'})
\sum_{\text{channels}}
\hat{\varphi}\bigl(\kappa_{x} - \kappa_{y}\bigr)\,
s(\kappa_{x}, \kappa_{y}),
\qquad
s(\kappa_{x}, \kappa_{y}) =
\tfrac{1}{2}\bigl(1 + e^{\,i(\theta_{x} - \theta_{y})}\bigr),
\]
the channels running over the $\sigma$-orbits
$\{\kappa, -\kappa\}$ of the two labels, $\theta$ the momentum
angle, and $c$ a non-vanishing kinematic factor. The spinor factor
$s$ vanishes precisely when the two momenta are antiparallel; a
$\sigma$-even $\varphi$ supported on a single momentum difference
isolates a conjugate channel pair, and the element is non-zero
whenever the isolated pair does not cancel identically.
\end{lemma}

\begin{proof}
Plane-wave eigenspinors of the flat operator have positive-energy
polarisations $u_{\kappa} \propto (1, e^{i\theta_{\kappa}})$; the
conformal perturbation is multiplicative up to the kinematic
dressing of the covariance formula, so its element between plane
waves is the Fourier coefficient of $\varphi$ at the momentum
transfer times the polarisation overlap
$\bar{u}_{\kappa_{x}} u_{\kappa_{y}}$, which is the displayed $s$;
$\sigma$-projection replaces each label by its two-element orbit
and the element by the corresponding finite sum, the orbits being
free because $c_{1}(L) = 3$ is odd \cite{MEF3}. The vanishing
condition $\theta_{x} - \theta_{y} = \pi$ is antiparallelism.
\end{proof}

\begin{proposition}[Non-degeneracy of the gating at the lowest
tower crossings]\label{ndc:verdict}
For each base tower, the crossing loci of the lowest tower level
with the harmonic block exist over an interval of base eigenvalues
$\lambda_{b}^{2}$ --- the deformation tunes the fibre gap through
the coincidence, so no exact matching of $\lambda_{b}^{2}$ to a
reference gap is required --- and at every such crossing among the
lowest pairs of the reference spectrum the gating element
\eqref{eq:gating} is non-zero for an explicit admissible product
deformation: the base factor by Lemma~\textup{\ref{ndc:base}}, the
fibre factor because none of the lowest pairs is antiparallel and
no channel cancellation occurs among them. We refer to this
statement as NDC. The gating amplitudes are bounded away from zero
on the verified list.
\end{proposition}

\begin{proof}
The fibre gaps of the reference spectrum form the integer list
\eqref{eq:levellist}; the list, its multiplicities against the
reference table, the gap values, and the non-antiparallelism of the
ten lowest crossing pairs are machine-verified (certificate
\texttt{NDC\_certificate.lean}). The channel factors of the ten
lowest pairs are bounded below by $0.80$ in the dominant channel,
so no cancellation occurs. Interval rather than point existence of
the crossing follows from the first-order movability of the fibre
levels under mean-free conformal deformations, the derivative being
the eigenspinor density pairing of the conformal variation formula
\cite{BourguignonGauduchon1992}.
\end{proof}

\begin{remark}[What is already forced]\label{rb1:free}
Independently of everything in the next section, under the smoothed
window the strata of Definition~\textup{\ref{rb1:strata}} have
codimensions three and above; hence the admissible region is
connected and simply connected, $b_{0} = 1$ and $b_{1} = 0$, and
since the links of the generic strata --- spheres and products of
spheres --- carry even cohomology only, every odd Betti number of
the complement vanishes. The evenness and lacunarity of the target
sequence, the properties on which the Morse argument of
Section~\ref{sec:count} runs, are structural consequences of the
symmetry class of the operator family, not accidents of the
geometry.
\end{remark}

\section{The collapse and the ring}\label{sec:collapse}

The purpose of this section is to convert the boundary structure of
the in-window strata into linear relations among the linking
classes, to collapse each base-content sector to a single
generator, and to fix the ring relations by a transgression pairing
against the base intersection forms. Three hypotheses enter here
and are displayed in full; each is stated once, with what it
suppresses named.

\begin{lemma}[Edge capping and the increment
relation]\label{rb2:edge}
Let $\Sigma_{c}^{w}$ be a charged in-window component along which
$dx_{c} \neq 0$. Near an edge point the excised set is a $C^{1}$
submanifold-with-boundary of codimension three, and a linking
two-sphere $S_{c}$ transported along an admissible ambient path on
which $x_{c}$ exits the window is capped by a three-chain in the
complement of $\Sigma_{c}^{w}$. The cap meets the remaining
in-window components transversally in finitely many points, and in
$H_{2}(\Xad(\varepsilon); \Q)$
\begin{equation}\label{eq:increment}
[S_{c}] \;=\; \sum_{c' \neq c} n_{cc'}\,[S_{c'}],
\end{equation}
the $n_{cc'}$ signed intersection numbers of the cap with
$\Sigma_{c'}^{w}$. In particular a component whose window exit is
realisable by a path whose cap meets no other charged component has
trivial linking class.
\end{lemma}

\begin{proof}
Transversality of $x_{c}$ to the endpoint value makes
$\overline{\Sigma_{c}^{w}}$ locally a product of the
codimension-three normal model with a half-line in the $x_{c}$
direction; the linking sphere in a normal slice on the in-window
side is slid past the boundary through the normal model of the
boundary stratum, where the excised set is absent, and contracted
in the ambient contractible space
(Proposition~\ref{rb:ambient}). The resulting three-chain is put in
general position with respect to the remaining strata: a
three-chain meets codimension-three strata in isolated points and
misses codimension-six strata, the dimension arithmetic being
machine-verified (certificate \texttt{RB2\_certificate.lean}). Each
isolated intersection contributes the linking class of the stratum
met, with the sign of the local intersection number, which is the
displayed relation. The final clause is an Alexander-type
triviality: a half-infinite codimension-three sheet in a
contractible ambient space links nothing, and only the transversal
meetings with other charged strata obstruct the capping.
\end{proof}

\begin{lemma}[Sector separation]\label{rb2:sector}
The relations of Lemma~\textup{\ref{rb2:edge}} connect components
within a single base-content sector only: fibre components to fibre
components, $\Stwo$-tower to $\Stwo$-tower, $\CPtwo$-tower to
$\CPtwo$-tower. An exchange between components of distinct sectors
would require the capping chain to meet the simultaneous
coincidence of two independently gated pairs, a stratum of
codimension six, which a generic three-chain misses; and the
bidegree labels of the sector classes are rigid under admissible
deformation, the base being $\sigma$-inert and undeformed.
\end{lemma}

\begin{proof}
Two coincidence conditions on pairs with distinct base content
impose independent complex Hermitian degeneracies, three real
conditions each by Lemma~\ref{rb1:gate}; the intersection has
codimension six and the dimension count $3 - 6 < 0$ is
machine-verified. Rigidity of the bidegree labels is
Lemma~\ref{o3:blocks} below, whose proof is independent of this
section: the residual $\Ztwo$ is a fibre isometry and acts
trivially on base cohomology.
\end{proof}

The remaining input of the collapse is a Gram non-degeneracy: for
each pair of consecutive same-sector in-window components, the Gram
matrix of the eigenspinor densities governing the first variation
of the two crossing positions must be non-degenerate on the
admissible deformations, so that a one-parameter admissible family
realises the exit of the first and the entry of the second as a
single transversal exchange, the capping chain of
Lemma~\ref{rb2:edge} meeting the entering component with
intersection number $\pm 1$. We refer to this statement as GD. Like
NDC it is a finite check against the level list
\eqref{eq:levellist}, and like NDC it is established at reduction
level on the verified list; the governing densities are the
transition densities of the crossing pair --- the diagonal
densities are exactly constant at the flat configuration by the
density flattening of \cite{MEF3} and pair to zero with every
mean-free direction --- and through the conformal variation formula
\cite{Kato, Hitchin1974, Hijazi, BourguignonGauduchon1992} the
movability functional of a component is represented in $L^{2}$ by
an element supported on the pair's non-antiparallel momentum
transfers, with coefficient the polarisation overlap of
Lemma~\textup{\ref{ndc:fibre}}. Distinct canonical frequencies are
orthogonal, so the Gram matrix of two components is degenerate only
if their non-vanishing Fourier supports coincide as sets, and GD
reduces to support arithmetic. For the tower sectors the coupling
direction is supplied by the gating integral, so GD there is
consumed together with NDC.

\begin{proposition}[GD on the verified list]\label{gd:verdict}
On the reference spectrum, in the integer momentum model
$K(m, j) = (100m,\, 93j)$, $j = 2n - 3$: the fibre-sector
components are carried by the seven multiplicity-four levels among
the ten lowest, each with transfer support
$\{(0, 186j),\, (200m, 0)\}$, both channels non-antiparallel since
$mj \neq 0$, and the seven supports pairwise distinct; and the
tower-sector components of the ten crossing pairs of
Proposition~\textup{\ref{ndc:verdict}} have pairwise distinct,
non-empty transfer supports, with no vanishing polarisation
coefficient on any supporting channel and the normalised Gram
determinants of consecutive components bounded below by $0.49$
under worst-case phases. Consequently every pair of same-sector
components on the list --- in particular every adjacent pair, under
any ordering convention for the sector chains --- has
non-degenerate Gram matrix: GD holds on the verified list.
\end{proposition}

\begin{proof}
The support computations, the pairwise set-distinctness (all
twenty-one fibre pairs and all forty-five tower pairs, strictly
stronger than the adjacent pairs the collapse consumes), the
non-antiparallelism of every supporting channel, and the
non-emptiness of every support are machine-verified in the integer
model (certificate \texttt{GD\_certificate.lean}); the polarisation
weights and the Gram lower bounds are numerical evaluations
recorded per pair. Independence follows from Fourier orthogonality:
distinct supports with non-vanishing coefficients exclude
proportionality, and the worst-case-phase bound controls the
determinant away from zero. The residuals are those of
Proposition~\ref{ndc:verdict} and of the same kind: the standard
first-order perturbation input at the degenerate block, and the
recorded numerical bounds.
\end{proof}

\begin{hypothesis}[OW --- window retention and
detection]\label{rb2:ow}
For each sector the curvature two-form of the sector's spectral
pair bundle, smoothly windowed, extends to a closed form on
$\Xad(\varepsilon)$ representing a non-zero class dual to the
sector's linking spheres; and the three sector classes so detected
are linearly independent in $H^{2}(\Xad(\varepsilon); \Q)$.
Equivalently, no admissible path drains a sector's charged
in-window set while capping its last linking sphere trivially. The
deflationary reading is stated at once: the collapse machinery
below proves rank at most one per sector; that the rank is exactly
one, three times over and independently, is OW and nothing else.
The hypothesis is the single home of every lower bound in this
paper.
\end{hypothesis}

\begin{proposition}[The collapse: one class per
sector]\label{rb2:collapse}
GD and NDC hold at reduction level on the verified list
\textup{(Propositions~\ref{gd:verdict} and~\ref{ndc:verdict})}.
Then within each sector all in-window linking classes are
rationally proportional: the increment lattice of the sector has
rank at most one, and any chain of adjacent exchanges identifies
the classes up to sign. Assuming
Hypothesis~\textup{\ref{rb2:ow}}, each sector contributes exactly
one generator --- $z$ for the fibre sector, $y$ for the $\Stwo$
tower, $x$ for the $\CPtwo$ tower --- and the three are
independent.
\end{proposition}

\begin{proof}
By Proposition~\ref{gd:verdict} each adjacent pair $(c, c')$ admits
the single-exchange family, and Lemma~\ref{rb2:edge} reads
$[S_{c}] = \pm[S_{c'}]$; Lemma~\ref{rb2:sector} confines every such
relation to the sector. Consecutive relations compose along the
energy-ordered chain of the sector's in-window components, and the
sign bookkeeping of a $\pm 1$ chain --- every class equal to plus
or minus the first --- is machine-verified (certificate
\texttt{RB2\_certificate.lean}). Rank at most one follows;
Hypothesis~\ref{rb2:ow} supplies rank exactly one per sector and
the independence of the three generators.
\end{proof}

\begin{hypothesis}[TR --- transgression
compatibility]\label{rb2:tr}
The tower generators are detected by evaluation against base
two-cycles: the family of admissible profiles obtained by
transporting a tower-crossing configuration over a two-cycle of $B$
through the gated coupling functional defines a pairing under which
the product structure of $H^{*}(\Xad(\varepsilon); \Q)$ on the
tower generators intertwines the intersection ring of the
corresponding base factor \cite{BottTu1982}. The hypothesis is the
single mapping-space statement of the paper; it suppresses nothing
about existence or rank, which belong to OW, and carries exactly
the ring relation of the next proposition.
\end{hypothesis}

\begin{proposition}[The ring discriminator]\label{rb2:ring}
Assume Hypothesis~\textup{\ref{rb2:tr}}. Then $x^{2} \neq 0$ and
$x^{3} = 0$, while $y^{2} = 0$: the truncated relation of the
$\CPtwo$ generator and the exterior relation of the $\Stwo$
generator are forced by the base intersection forms
$H^{*}(\CPtwo; \Q) = \Q[a]/(a^{3})$ and
$H^{*}(\Stwo; \Q) = \Q[b]/(b^{2})$, and the asymmetry between the
two tower generators is a statement of base topology alone:
whichever complement class has surviving square is necessarily the
$\CPtwo$-tower class.
\end{proposition}

\begin{proof}
Under Hypothesis~\ref{rb2:tr} the square of a tower generator pairs
with the self-intersection of the base two-cycle it is evaluated
against: for the $\CPtwo$ tower the line class satisfies
$[\mathbb{CP}^{1}] \cdot [\mathbb{CP}^{1}] = [\mathrm{pt}]$, non-zero in
$H^{4}(\CPtwo; \Q)$, and the cube meets $a^{3} = 0$; for the
$\Stwo$ tower $H^{4}(\Stwo; \Q) = 0$ and the square vanishes. The
fibre generator $z$ carries $z^{2} = 0$ as the linking class of a
smooth component. The exponent arithmetic of the truncated ring,
and the discriminating computation $a^{2} \neq 0$, $a^{3} = 0$,
$b^{2} = 0$, are machine-verified (certificate
\texttt{RB2\_certificate.lean}).
\end{proof}

\begin{hypothesis}[HS --- dependence of the higher
strata]\label{hyp:hs}
The linking classes of the transverse intersection strata --- the
residual strata of Definition~\textup{\ref{rb1:strata}} in which
two or more independently gated coincidences occur simultaneously
--- are products of the sector generators, contributing no new
generators to $H^{*}(\Xad(\varepsilon); \Q)$. The hypothesis
suppresses an exotic class borne by a deep stratum; for transverse
configurations the links are products of the links of the factors
and the statement holds generically, so what is suppressed is a
failure of genericity at a multiple coincidence.
\end{hypothesis}

\section{Stabilisation and the generator-level
theorem}\label{sec:stab}

The exchange relation that collapses a sector is the same
transversal meeting that exhausts the defect of the retraction at
finite energy: the collapses and the stabilisation are one
bookkeeping run in opposite directions. This section makes that
statement precise and assembles Theorem~\ref{thm:A}.

\begin{theorem}[Stabilisation]\label{rb2:stab}
Assume Hypotheses \textup{OW} and \textup{HS}; GD and NDC hold at
reduction level on the verified list
\textup{(Propositions~\ref{gd:verdict} and~\ref{ndc:verdict})}.
Then the rational
homology of the truncations $\Xad(\varepsilon)$ stabilises at a
finite $\varepsilon^{*}$: for
$\varepsilon' \geq \varepsilon \geq \varepsilon^{*}$ the inclusions
$\Xad(\varepsilon) \hookrightarrow \Xad(\varepsilon')$ induce
isomorphisms on rational homology, and $\varepsilon^{*}$ may be
taken as the smallest $\varepsilon$ for which
$\varepsilon^{2}\gamma^{2}$ exceeds the base Dirichlet energy of
the lowest-harmonic normal deformations opening the lowest crossing
of each tower --- explicit in the first non-zero base eigenvalues
$\lambda_{1}(\CPtwo)$, $\lambda_{1}(\Stwo)$ and the gating
amplitudes of Proposition~\textup{\ref{ndc:verdict}}. Consequently,
for every $\varepsilon \geq \varepsilon^{*}$ the inclusion
$\Xad(\varepsilon) \hookrightarrow X_{B}$ is a rational homology
isomorphism, and the truncated and unrestricted formulations of
every statement in this paper are one statement.
\end{theorem}

\begin{proof}
By Lemma~\ref{rb:colimit} every class of the full space is realised
at some finite $\varepsilon$. Pushing a $k$-cycle down the
retraction of Definition~\ref{rb:retraction} sweeps a
$(k+1)$-trace, which meets the codimension-three crossing loci, for
$k = 2$, in isolated points and misses the codimension-six strata,
the dimension arithmetic machine-verified; each meeting changes the
class by a linking-sphere increment, so every defect of the
retraction between truncations is a sum of linking increments of
charged codimension-three strata. By
Proposition~\ref{rb2:collapse} every such increment is rationally
proportional to its sector generator, and the generator is realised
on the linking sphere of the sector's lowest in-window crossing:
for the fibre sector this sphere lies over the base-constant slice
and is present in every truncation; for each tower sector its
normal directions require the base-dependent gating perturbation,
which is opened by the lowest tower harmonic at the base Dirichlet
cost displayed, hence contained in $\Xad(\varepsilon)$ once
$\varepsilon \geq \varepsilon^{*}$. Higher strata contribute only
products of the generators by Hypothesis~\ref{hyp:hs}. For
$\varepsilon' \geq \varepsilon \geq \varepsilon^{*}$ the inclusion
therefore hits every generator and every product, and induces an
isomorphism on rational homology; the final clause is the colimit
of a system whose maps are isomorphisms beyond $\varepsilon^{*}$,
attained at $\varepsilon^{*}$.
\end{proof}

\begin{proof}[Proof of Theorem~\textup{\ref{thm:A}}]
Existence and one-per-sector of the generators is
Proposition~\ref{rb2:collapse} with Lemma~\ref{rb1:gate} and
Proposition~\ref{ndc:verdict}; the ring relations are
Proposition~\ref{rb2:ring} for $x$ and smooth-component linking for
$y$ and $z$; the higher strata contribute the products and nothing
else by Hypothesis~\ref{hyp:hs}. The cohomology is therefore the
truncated polynomial ring displayed, and its Betti sequence is the
coefficient list of $(1 + t^{2} + t^{4})(1 + t^{2})(1 + t^{2})$,
namely $(1, 0, 3, 0, 4, 0, 3, 0, 1)$ --- the ring and the sequence
of $H^{*}(\Keight; \Q)$ itself, the identification that carries the
remainder of the paper --- the monomial count being
machine-verified twice (Proposition~\ref{prop:betti}).
Stabilisation and the equivalence of the truncated and unrestricted
formulations are Theorem~\ref{rb2:stab} over
Lemma~\ref{rb:colimit}; the ambient spaces carry no topology
(Proposition~\ref{rb:ambient}), so every class is
discriminant-borne, as claimed.
\end{proof}

\begin{remark}[Homology suffices]\label{lp0b:homology}
Theorem~\ref{thm:A} is a rational homology statement at generator
level, not a homotopy equivalence. The Morse argument of
Section~\ref{sec:count} consumes exactly the Betti numbers of the
space the functional lives on, and Theorem~\ref{thm:A} supplies
them; nothing below asks for more.
\end{remark}

\begin{remark}[The open set, displayed]\label{lp0b:open}
The hypotheses of Theorem~\ref{thm:A} divide by kind. The two
finite non-degeneracy checks against the level list
\eqref{eq:levellist} are both established at reduction level on
the verified list: NDC with the channel factors bounded below
(Proposition~\ref{ndc:verdict}), GD with pairwise-distinct
transfer supports and the Gram determinants bounded below
(Proposition~\ref{gd:verdict}). OW and TR are the two named
structural residuals: OW
carries every lower bound (that the three classes exist and are
independent), TR carries the ring relation of the $\CPtwo$
generator. HS is the dependence of the higher strata, generic for
transverse configurations. Nothing else is open on the statement.
\end{remark}

\section{The critical count}\label{sec:count}

The purpose of this section is to convert the topology of
Theorem~\ref{thm:A} into a count of critical points. Three
ingredients enter: the six base-sector replicas of the constant
stratum must be six distinct critical strata rather than one
stratum with internal degeneracy (Section~\ref{sec:six}); the fibre
counting factor, which takes the value two at the pillowcase under
three a priori different descriptions, must be identified
(Section~\ref{sec:two}); and the Morse inequalities must be
consumed (Section~\ref{sec:morse}).

\subsection{The six replicas are six}\label{sec:six}

\begin{lemma}[Second variation at the constant
stratum]\label{o3:secondvar}
At a constant configuration, with no level on a window edge, the
second variation of $\omega$ in slice-admissible directions
$\delta\Phi$ is a well-defined quadratic form $Q[\delta\Phi]$,
assembled from the second-order Kato expansion of the eigenvalues
of $D_{\Phi}$ through the conformal variation formula
\cite{Kato, BourguignonGauduchon1992}, the block-trace contraction
at degenerate levels, and the second derivative of the Dyson--Mehta
functional in the unfolded positions; each contribution is dressed
by the base multiplicity of the level at which it enters.
\end{lemma}

\begin{lemma}[Bidegree block decomposition]\label{o3:blocks}
At the adiabatic point the spectrum carries the block form
$\lambda^{2} = \lambda_{b}^{2} + \lambda_{f}^{2}(\Phi)$, and the
six base harmonic states of $B = \CPtwo \times \Stwo$ are labelled
by the pairwise-distinct bidegrees
\[
(p, q) \;\in\; \{(0,0),\, (2,0),\, (0,2),\, (4,0),\, (2,2),\,
(4,2)\},
\qquad p \in H^{p}(\CPtwo),\; q \in H^{q}(\Stwo),
\]
the coefficient list of $(1 + t^{2} + t^{4})(1 + t^{2})$ being
$(1, 0, 2, 0, 2, 0, 1)$ with total $6 = \chi(B)$. The quadratic
form of Lemma~\textup{\ref{o3:secondvar}} accordingly decomposes
into six blocks indexed by these bidegrees, and the bidegree labels
are preserved by the residual symmetry of the problem, which acts
on the fibre alone. The label arithmetic --- six blocks, pairwise
distinct, exhausting the products $\{0, 2, 4\} \times \{0, 2\}$ ---
is machine-verified (certificate \texttt{O3\_certificate.lean}).
\end{lemma}

\begin{proof}
The block form at the adiabatic point and the multiplicity six are
the replication mechanism of the assembly proposition of
\cite{MEF3}: at the adiabatic point the configuration varies only
the fibre part and every fibre level's variation is replicated once
per base state, the unsigned harmonic multiplicity of the base
equalling $\chi(B) = 6$ because the Hodge diamond of $B$ is purely
diagonal. The bidegree labelling is K\"unneth bookkeeping; label
preservation holds because the residual $\Ztwo$ is a fibre isometry
and acts trivially on base cohomology.
\end{proof}

\begin{lemma}[Factor-scaling response of the harmonic
dressings]\label{o3:scaling}
Let
$g(a_{1}, a_{2}) = e^{2a_{1}}g_{\CPtwo} \oplus e^{2a_{2}}g_{\Stwo}$
be the two-parameter family of factor-wise conformal scalings of
the base metric, an admissible direction of the identified
conformal and warp channel. A harmonic representative of bidegree
$(p, q)$ has pointwise norm scaling as
\[
\bigl|\alpha_{(p,q)}\bigr|^{2}_{g(a_{1}, a_{2})}
\;=\; e^{-2p\,a_{1} - 2q\,a_{2}}\,
\bigl|\alpha_{(p,q)}\bigr|^{2}_{g},
\]
the exponent weighted by the factor degrees. Consequently the
dressed weight functional of the bidegree-$(p, q)$ block of
Lemma~\textup{\ref{o3:blocks}} responds to the scaling channel with
logarithmic gradient $-2(p, q)$ at the constant stratum.
\end{lemma}

\begin{proof}
The pointwise norm of a $p$-form on a factor rescaled by $e^{2a}$
carries $p$ inverse-metric contractions, each contributing
$e^{-2a}$; the factors multiply across the product base. That the
degree-weighted exponent is the channel by which the base state
dresses the balance condition is the replication mechanism of
Lemma~\ref{o3:blocks} read on the identified conformal and warp
channel of Definition~\ref{rb:space}: the dressing of each block is
carried by the harmonic density of its base state, and the harmonic
states track their bidegree representatives under the
correspondence of Lemma~\ref{o3:blocks}.
\end{proof}

\begin{proposition}[Pairwise inequality of the six dressed
weights]\label{o3:pairwise}
The six response vectors of Lemma~\textup{\ref{o3:scaling}} are the
six bidegrees
$(0,0),\, (2,0),\, (0,2),\, (4,0),\, (2,2),\, (4,2)$, which are
pairwise distinct; in particular the two pairs of equal total
degree, $(2,0) \neq (0,2)$ and $(4,0) \neq (2,2)$, are separated by
the anisotropy of the family, which a single overall scaling cannot
achieve. For every pair of blocks there is therefore an explicit
admissible direction on which the two dressed weight functionals
respond with different derivatives, so the six functionals are
pairwise unequal; and by the label rigidity of
Lemma~\textup{\ref{o3:blocks}} no admissible deformation can merge
two blocks whose weights differ. The six base-sector replicas of
the constant stratum are six distinct critical strata, and the
internal-degeneracy alternative is excluded. The distinctness and
the separation of the equal-degree pairs are machine-verified
(certificate \texttt{O3\_pairwise\_certificate.lean}).
\end{proposition}

\begin{proof}
Distinctness of the six vectors is finite arithmetic. Two
functionals with different derivatives along one admissible
direction are unequal as functionals; one separating direction per
pair is supplied by the response vectors, the difference of any two
being a non-zero vector realised by an anisotropic scaling.
Exclusion of merging then follows: a merged stratum would require
an operator identity equating dressed weights across distinct
bidegrees --- a strictly stronger event than degeneracy of
eigenvalues --- contradicting the exhibited inequality.
\end{proof}

The remaining obligation on this argument for an unconditional
grade is one finite check: the warp-elimination identities are
exact at first order on the ground sector, and their second-order
use in Lemma~\ref{o3:secondvar} is to be verified against them
once, not assumed. We record this in
Section~\ref{sec:outlook}.

\subsection{The fibre counting object}\label{sec:two}

At the pillowcase three a priori different quantities take the
value two. No result in this paper consumes the identification of
which: Theorem~\ref{thm:B} counts through the Betti numbers and
Hypothesis~\textup{UB}, and the factor two enters only the
structural description of the critical set. The purpose of this
subsection is to identify that factor, not to feed the count. The
three must first be separated
before the counting object can be named, and this subsection
constructs the family on which they separate.

\begin{definition}[The three candidate counting
objects]\label{o4:objects}
For the pillowcase fibre with its decoration and the residual
$\Ztwo$ of the symmetric-criticality theorem, set
\[
N_{\chi} = \chi_{\mathrm{top}}(\Ttwo/\Ztwo),
\qquad
N_{\mathrm{stab}} = \#\{\text{stabiliser classes of the
$\Ztwo$-action on } \Ttwo\},
\qquad
N_{\mathrm{corner}} = \#\{\text{corner orbits under the residual
swap}\},
\]
the fibre Euler characteristic, the stabiliser-class count, and the
corner-orbit count respectively.
\end{definition}

\begin{lemma}[Coincidence at the pillowcase]\label{o4:coincide}
$N_{\chi} = N_{\mathrm{stab}} = N_{\mathrm{corner}} = 2$, each by
an independent computation: the invariant rational cohomology of
the pillowcase is $(1, 0, 1)$, giving $N_{\chi} = 2$; the
$\Ztwo$-action on $\Ttwo$ has exactly the two stabiliser classes,
trivial and full, giving $N_{\mathrm{stab}} = 2$; and the four
corners fall into two pairs under the residual swap, which is free
on them, giving $N_{\mathrm{corner}} = 2$. The three computations
are machine-verified independently (certificate
\texttt{O4\_certificate.lean}). The coincidence at $2$ is the
problem, not the result.
\end{lemma}

\begin{definition}[The group-replacement family]\label{o4:family}
For $n \in \{3, 4, 6\}$ let $\Ttwo_{n}$ be the torus at the rigid
modulus admitting the order-$n$ rotation $\rho$
($\tau = e^{i\pi/3}$ for $n = 3, 6$; $\tau = i$ for $n = 4$), $R$
the integer matrix of $\rho$ on a lattice basis, and
$\Ttwo_{n}/\Zn$ the quotient orbifold. The family deforms the fibre
away from the pillowcase modulus; this is permitted by the framing
of the separation test, which requires only that the counts and the
critical structure remain defined on each member, and the modulus
displacement is recorded rather than suppressed.
\end{definition}

\begin{lemma}[Equivariant $\Spinc$ classes]\label{o4:equivariance}
The rotation acts on the four $\Spinc$ shift classes
$\tfrac{1}{2}\Lambda^{*}/\Lambda^{*} \cong (\Ztwo)^{2}$ through
$R \bmod 2$. The invariant classes are: all four at $n = 2$; the
trivial class only at $n = 3$ and $n = 6$; the trivial class and
$(\tfrac{1}{2}, \tfrac{1}{2})$ at $n = 4$. Consequently the
non-integral shift --- the carrier of the freeness mechanism ---
extends $\rho$-equivariantly to the $n = 4$ member alone; at
$n = 3, 6$ the equivariant problem forces the integral class, the
rotation fixes the zero label, and the freeness mechanism of the
odd-degree protection does not extend. The class arithmetic is
machine-verified (certificate
\texttt{O4\_family\_certificate.lean}).
\end{lemma}

\begin{proof}
Invariance of a class $v \in (\Ztwo)^{2}$ is
$Rv \equiv v \bmod 2$, a finite check per member. Freeness of the
rotation on the shifted lattice $k_{0} + \Lambda^{*}$ requires
$0 \notin k_{0} + \Lambda^{*}$, since $1 - \rho^{j}$ is invertible
for $0 < j < n$; the integral class violates this at the zero
label, and the non-integral classes restore it wherever they
extend.
\end{proof}

\begin{lemma}[Residual translations and their signs on the
family]\label{o4:residual-n}
The translations descending to symmetries of $\Ttwo_{n}/\Zn$ form
the group $\Lambda/(1 - \rho)\Lambda$, of order
$|\det(1 - R)| = 4, 3, 2, 1$ for $n = 2, 3, 4, 6$. On the shifted
momentum lattice the conjugation of the rotation lift by a lifted
translation $T_{a}$ carries the constant phase
\[
\varepsilon_{a} \;=\;
e^{-i\langle k_{0},\, (1 - \rho)a\rangle},
\]
well-defined because $(1 - \rho)a \in \Lambda$ pairs integrally
with $\Lambda^{*}$; a translation survives to the projected problem
precisely when $\varepsilon_{a} = 1$. At $n = 2$ with the
pillowcase class this recovers the conjugation computation of
\cite{MEF3} exactly: the unshifted half-period survives and the two
shifted half-periods carry $\varepsilon = -1$. At $n = 4$ with the
equivariant class $(\tfrac{1}{2}, \tfrac{1}{2})$ the single
non-trivial candidate has
$\langle k_{0}, (1 - \rho)a\rangle = \pi$ and does not survive. At
$n = 3$ with the forced integral class all three candidates carry
$\varepsilon = 1$ and survive, acting simply transitively on the
three cone points. At $n = 6$ the candidate group is already
trivial.
\end{lemma}

\begin{proof}
$t_{a}\rho t_{a}^{-1} = t_{(1-\rho)a}\,\rho$ descends iff
$(1 - \rho)a \in \Lambda$, giving the group and its order as the
index $|\det(1 - R)|$. On a momentum state the conjugation phase is
$e^{-i\langle(1 - \rho^{-1})k,\, a\rangle}$; transposing the
orthogonal rotation moves $(1 - \rho^{-1})$ onto $a$ as
$(1 - \rho)$, and the $k$-dependence drops modulo $2\pi$ because
$(1 - \rho)a \in \Lambda$, leaving the shift term displayed. The
four sign evaluations are finite arithmetic on half-integral
classes, machine-verified; the $n = 2$ case is the recorded
consistency check against the source computation, which the general
formula reproduces verbatim before it is trusted on the new
members.
\end{proof}

\begin{proposition}[The separation table]\label{o4:table}
On the family the three candidate counting objects take the values
\[
\begin{array}{c|ccc}
n & N_{\chi} & N_{\mathrm{stab}} & N_{\mathrm{corner}} \\ \hline
3 & 2 & 2 & 1 \\
4 & 2 & 3 & 3 \\
6 & 2 & 4 & 3
\end{array}
\]
by the following routes. $N_{\chi} = 2$ on every member: the
rotation has no invariant vector on $H^{1}(\Ttwo; \Q)$ since
$\det(1 - R) \neq 0$, so the invariant Betti numbers are
$(1, 0, 1)$ throughout. $N_{\mathrm{stab}} = 2, 3, 4$: the
stabiliser classes occurring are read from the fixed-point counts
$|\mathrm{Fix}(\rho^{j})| = |\det(1 - R^{j})|$ --- at $n = 3$ the
classes $\{1, \Zn[3]\}$; at $n = 4$ the classes
$\{1, \Ztwo, \Zn[4]\}$, the two extra order-two points forming one
quotient cone point; at $n = 6$ the classes
$\{1, \Ztwo, \Zn[3], \Zn[6]\}$. $N_{\mathrm{corner}}$: each member
has three cone points; at $n = 3$ the three surviving translations
of Lemma~\textup{\ref{o4:residual-n}} identify them into one orbit;
at $n = 4$ and $n = 6$ no translation survives and the three points
are three orbits. Every numerical ingredient --- orders,
determinants, fixed-point counts, invariant classes, signs, orbit
arithmetic, and the three triples with their distinctness patterns
--- is machine-verified. In particular the coincidence of the
pillowcase breaks on every member, and at $n = 6$ the three
quantities are pairwise distinct.
\end{proposition}

The identification now reduces to two readouts: the $n = 4$
member, where the equivariant non-integral class preserves the
freeness mechanism, decides between $N_{\chi} = 2$ and
$N_{\mathrm{stab}} = N_{\mathrm{corner}} = 3$; the $n = 6$
member, on which all three candidates are pairwise distinct,
decides independently. We execute both. Two structural lemmas
prepare the ground.

\begin{lemma}[Flattening is exclusive to the
involution]\label{o4r:flat}
Let $\chi$ be a character-projected eigenspinor built on a free
rotation orbit of the $\Zn$ member. The density $|\chi|^{2}$
carries cross-terms between the orbit members $\rho^{j'}\kappa$,
$\rho^{j}\kappa$ with polarisation factor proportional to
$1 + e^{i(j - j')\alpha}$, $\alpha = 2\pi/n$, which vanishes
exactly when $(j - j')\alpha = \pi$: on the antipodal terms
alone. At $n = 2$ every cross-term is antipodal and the density is
exactly constant --- the flattening mechanism protecting the
constant stratum in \cite{MEF3}. At $n = 4$ and $n = 6$ the
antipodal terms vanish (since $\rho^{n/2} = -1$) while the
remaining cross-terms survive with non-zero exact coefficients; at
$n = 3$ no cross-term is antipodal and none vanishes. The
character-projected densities are therefore non-constant on every
member other than the pillowcase: the pointwise flattening is a
joint peculiarity of the antipodal involution and the odd
decoration, and the critical structure on the members can be
carried only by the homological machinery, which is
group-agnostic.
\end{lemma}

\begin{proof}
The rotation lift carries eigenspinors to eigenspinors with a
constant phase, so the cross-term amplitude between orbit members
is the polarisation overlap of Lemma~\ref{ndc:fibre} at angle
difference $(j - j')\alpha$. The vanishing pattern ---
$1 + i^{2} = 0$ with $1 + i^{\pm 1} \neq 0$ in the Gaussian
integers for $n = 4$; $1 + z^{3} = 0$ with $1 + z^{j} \neq 0$,
$j = 1, 2, 4, 5$, for $z$ a primitive sixth root for $n = 6$, the
$j = 2, 4$ cases covering $n = 3$ --- is exact ring arithmetic,
machine-verified together with $\rho^{2} = -1$ and
$\rho^{3} = -1$ on the respective lattices (certificate
\texttt{O4\_readout\_certificate.lean}).
\end{proof}

\begin{lemma}[The fibre-sector structure of the
members]\label{o4r:structure}
On the $n = 4$ member (doubled coordinates, momenta odd-odd,
$Q = a^{2} + b^{2}$) the six lowest shells are
$Q = 2, 10, 18, 26, 34, 50$ with counts $4, 8, 4, 8, 8, 12$, every
orbit of size four, so the character-projected multiplicities are
$1, 2, 1, 2, 2, 3$ and the recombination components sit at the
multi-orbit shells. On the $n = 6$ member (hexagonal norm
$Q = a^{2} + ab + b^{2}$; the zero label of the forced integral
class is the kernel and decouples from every first variation, the
factor $\lambda_{k}$ of the conformal variation formula vanishing
there) the eight lowest non-zero shells are
$Q = 1, 3, 4, 7, 9, 12, 13, 16$ with counts
$6, 6, 6, 12, 6, 6, 12, 6$, every orbit of size six, recombination
components at the count-twelve shells. On the $n = 3$ member the
orbits have size three, the antipode of every label lies outside
its orbit, and recombination pairs occur at every shell. In every
case the recombination components' transfer supports are pairwise
distinct and non-empty: the Gram non-degeneracy pattern of
Proposition~\ref{gd:verdict} transfers to the members verbatim,
by the same certificate.
\end{lemma}

\begin{proposition}[The readout: the counting object is the fibre
Euler characteristic]\label{o4r:verdict}
Transport the counting machinery of this paper member-wise at its
displayed conditionality: the ambient contractibility and the
discriminant enumeration are group-agnostic; the collapse runs on
the members' recombination components under
Lemma~\textup{\ref{o4r:structure}} and the member analogue of
Hypothesis~\textup{OW}; the Morse consumption runs under the
member analogues of Hypotheses~\textup{PS} and~\textup{UB}. Then
the critical count per base state on each member equals the total
invariant Betti number of the fibre,
$1 + b_{2}^{\mathrm{inv}} = 2$, and against the separation table:
at $n = 4$, count $2 = N_{\chi} \neq 3 = N_{\mathrm{stab}} =
N_{\mathrm{corner}}$; at $n = 6$, count $2 = N_{\chi} \neq 4 =
N_{\mathrm{stab}} \neq 3 = N_{\mathrm{corner}}$; at $n = 3$,
count $2 = N_{\chi} = N_{\mathrm{stab}} \neq 1 =
N_{\mathrm{corner}}$. The counting object is the fibre Euler
characteristic $N_{\chi}$, uniquely among the three candidates:
selected at $n = 4$, confirmed independently at $n = 6$, the
identification closing by selection.
\end{proposition}

\begin{proof}
The fibre factor per base state is one (the index-zero stratum)
plus the number of independent fibre-sector generators of the
member's admissible region. By Lemma~\ref{o4r:structure} the
collapse machinery of Section~\ref{sec:collapse} applies on each
member: the recombination components are codimension-three charged
crossings, their supports are pairwise distinct with no vanishing
coefficient, adjacent exchange families exist, and each member's
fibre sector collapses to rank at most one; the member retention
hypothesis supplies rank exactly one. The factor is therefore
$1 + 1 = 2$ on every member, matching the invariant cohomology
$(1, 0, 1)$ of the quotient sphere: $N_{\chi} = 2$ throughout.
The exclusions are arithmetic against the separation table,
machine-verified. The stabiliser count exceeds the total Betti
number at $n = 4$ and $n = 6$, so it can equal the critical count
only through excess critical points, excluded on the displayed
pattern. The corner-orbit count at $n = 3$ falls below the Morse
lower bound $c_{\mathrm{total}} \geq b_{\mathrm{total}} = 2$
\cite{Palais1963, Milnor1963} --- its exclusion consumes no
no-excess hypothesis at all --- and fails again at $n = 4$ and
$n = 6$ by $3 \neq 2$.
\end{proof}

\begin{remark}[Why the homological candidate, and the hint
resolved]\label{o4r:why}
The identification has a one-sentence explanation: of the three
candidates, only the Euler characteristic is a homological
invariant of the quotient, and the Morse count is homological. The
corner-orbit count tracks the residual translation group, whose
surviving order varies along the family ($3, 0, 0$ non-trivial
translations at $n = 3, 4, 6$) independently of the homology; the
stabiliser count tracks the isotropy lattice, which grows with $n$
while the invariant cohomology does not; the three coincide at the
pillowcase because the involution's residual data and its homology
both happen to equal two.
\end{remark}

\begin{lemma}[Retirement of the odd-degree family]\label{o4:odd}
For $L = \Oline{2k+1}$, $k \in \Z$, the conjugation sign of the
shifted half-periods is
$e^{i\pi(2n' - (2k+1))} = -1$ for every level index $n'$ and every
$k$, since $2n' - (2k+1)$ is odd; the residual symmetry, hence the
corner-orbit count, is unchanged across the odd degrees, and
$N_{\chi}$, $N_{\mathrm{stab}}$ are unchanged trivially on the
fixed orbifold. Variation of the decoration through the odd degrees
therefore separates nothing, and the group-replacement family is
the discriminating one. The parity is machine-verified.
\end{lemma}

The observation that once stood here as a suggestive hint --- that
the corner-orbit count is the only candidate whose defining data
enter the criticality mechanism directly, through the conjugation
signs of Lemma~\ref{o4:residual-n} --- is resolved in the negative
by Proposition~\ref{o4r:verdict}, and the resolution is sharpened
by Lemma~\ref{o4r:flat}: the mechanism those signs protect is
itself exclusive to the involution, so the corner data could never
have carried a count that must survive onto members where the
mechanism is absent.

\subsection{The Morse consumption}\label{sec:morse}

\begin{hypothesis}[PS --- Morse-theoretic
admissibility]\label{o2:ps}
The functional $\omega$ is $C^{2}$ on the admissible region $X$,
satisfies the Palais--Smale condition on the fixed-area slice, and
its critical points are non-degenerate. The hypothesis suppresses
three failure modes at once: loss of compactness along
Palais--Smale sequences approaching the discriminant, insufficient
regularity of the spectral functional at degenerate blocks, and
degenerate critical points, for which the counting statements below
would hold only in the Morse--Bott form.
\end{hypothesis}

\begin{proposition}[Existence at every index]\label{o2:exist}
Assume the hypotheses of Theorem~\textup{\ref{thm:A}} and
Hypothesis~\textup{\ref{o2:ps}}. Then the critical counts of
$\omega$ on $X$ satisfy
\[
c_{i} \;\geq\; b_{i}(\Keight)
\qquad (i = 0, \dots, 8),
\]
so that, beyond the constant stratum, critical points exist with
multiplicities at least $3, 4, 3, 1$ at indices $2, 4, 6, 8$: at
least eleven non-constant critical points, without any min-max
construction.
\end{proposition}

\begin{proof}
Under Hypothesis~\ref{o2:ps} the Morse inequalities hold for
$\omega$ on the Banach manifold $X$
\cite{Palais1963, Milnor1963}, graded against the Betti numbers of
$X$; by Theorem~\ref{thm:A} these are the numbers
\eqref{eq:betti}. The displayed multiplicities are its entries.
\end{proof}

\begin{hypothesis}[UB --- no excess critical
points]\label{hyp:ub}
The total number of critical points of $\omega$ on $X$ is at most
the total Betti number, twelve. The hypothesis suppresses critical
points invisible to the homology: saddles occurring in cancelling
pairs, and any critical point beyond those the topology forces. It
is the hardest open statement in this paper and no proof strategy
for it is claimed here; Section~\ref{sec:retreat} derives the
uniqueness consequence by a route that does not consume it.
\end{hypothesis}

\begin{proof}[Proof of Theorem~\textup{\ref{thm:B}}]
By Proposition~\ref{prop:betti} the odd Betti numbers vanish, so
the total Betti number equals the Euler characteristic, twelve. By
Proposition~\ref{o2:exist}, $c_{i} \geq b_{i}$ at every index; by
Hypothesis~\ref{hyp:ub}, $\sum_{i} c_{i} \leq \sum_{i} b_{i}$. The
inequalities admit no slack and collapse termwise, $c_{i} = b_{i}$
for all $i$; the termwise collapse from equal totals is
machine-verified (certificate \texttt{O2\_certificate.lean}). In
particular $c_{0} = b_{0} = 1$ and the remaining critical points
are $3 + 4 + 3 + 1$ of index $2, 4, 6, 8$. The structural content
of the twelve is that of Sections~\ref{sec:six} and~\ref{sec:two}:
six distinct base-sector strata dressed by the fibre counting
factor two, identified in Proposition~\ref{o4r:verdict} as the
fibre Euler characteristic.
\end{proof}

We state the deflationary reading of Theorem~\ref{thm:B} at once,
so it cannot be over-read. What is proved is a conditional
perfection: the displayed hypothesis set converts the Morse
inequalities into equalities. Existence of the constant critical
point is unconditional \cite{MEF3}; uniqueness of the stable
critical point is exactly as conditional as the displayed set, and
the set is not empty.

\section{The independent route}\label{sec:retreat}

The purpose of this section is to derive $c_{0} = 1$ from two local
hypotheses alone, consuming neither the topology of
Theorem~\ref{thm:A} nor the count of Theorem~\ref{thm:B}. The two
routes are independent insurance on the load-bearing claim: the
route of Sections~\ref{sec:ambient}--\ref{sec:count} is the
structural story, and delivers rigidity at every index; the present
route is two local statements, and delivers the uniqueness of the
stable critical point with no claim at higher index.

\begin{hypothesis}[B --- index cost of symmetry
breaking]\label{o1:break}
Every critical point of $\omega$ on $X \setminus X^{\Ztwo}$ has
Morse index at least one. Equivalently: at any critical point not
fixed by the residual $\Ztwo$, the second variation of $\omega$ is
negative in at least one direction. The hypothesis suppresses the
possibility of a symmetry-breaking local minimum; it does not
assert that non-symmetric critical points are absent, only that
none is stable. The natural attack is rigidity: symmetry breaking
costs $\omega$ at second order, the $\sigma$-partner pairing
forcing the breaking mode and its partner to contribute with
opposite signs, supplying the negative direction; this is recorded
as the attack route and not claimed.
\end{hypothesis}

\begin{hypothesis}[U --- symmetric uniqueness at index
zero]\label{o1:symuniq}
The constant stratum is the unique index-zero critical point of
$\omega$ restricted to $X^{\Ztwo}$. The hypothesis suppresses a
second, non-constant symmetric local minimum; it does not constrain
symmetric saddles. Two attack routes are recorded: geodesic
convexity of $\omega$ on $X^{\Ztwo}$ in a neighbourhood exhausting
the index-zero sublevel set; and triviality of the kernel of the
linearised balance equation \eqref{eq:el} at a putative second
minimum modulo the residual $\Ztwo$, with a continuation argument
along the slice connecting it to the constant stratum. The second
route would deliver uniqueness by elliptic rigidity, independent of
the Morse framework.
\end{hypothesis}

\begin{proof}[Proof of Theorem~\textup{\ref{thm:C}}]
Existence is unconditional: the constant configurations are
critical by the symmetry-protected criticality of the window theory
\cite{MEF3}, and lie in $X^{\Ztwo}$. For the count, let $p$ be any
critical point of index zero. By Hypothesis~\ref{o1:break},
$p \in X^{\Ztwo}$; the index of $\omega|_{X^{\Ztwo}}$ at $p$ is at
most the ambient index, hence zero, so $p$ is an index-zero
critical point of the restriction. By
Hypothesis~\ref{o1:symuniq} there is exactly one such, the constant
stratum. Hence $c_{0} = 1$. The counting step --- that a partition
of the critical set into a symmetric part with a single index-zero
member and a non-symmetric part with no index-zero member forces
$c_{0} = 1$ --- is machine-verified (certificate
\texttt{O1\_certificate.lean}).
\end{proof}

\begin{remark}[What the local route buys and what it
forgoes]\label{o1:scope}
Theorem~\ref{thm:C} consumes neither the configuration-space
topology nor the total count: the load is carried entirely by the
two displayed hypotheses, which are local and equivariant rather
than global and topological. What is forgone is perfection: without
the hypotheses of Theorem~\ref{thm:B} the equalities
$c_{i} = b_{i}$ at indices $i \geq 2$ are not claimed, and
non-symmetric saddles, if any exist, are simply not counted. The
statement that there is exactly one stable critical point survives
untouched. Corollary~\ref{cor:main} is immediate from
Theorems~\ref{thm:B} and~\ref{thm:C}.
\end{remark}

\section{Concluding remarks}\label{sec:outlook}

We close by recording, in one place, what is open and in what form;
the discipline throughout has been that every open statement is
displayed where it is consumed, so this section adds locations and
kinds, not new content.

The two finite non-degeneracy checks of the topology are closed at
reduction level on the verified list: NDC in
Proposition~\ref{ndc:verdict}, with the channel factors bounded
below, and GD in Proposition~\ref{gd:verdict}, with
pairwise-distinct transfer supports and the Gram determinants
bounded below, both machine-verified in the integer model; and the
critical-count readout on the group-replacement family is executed
in Proposition~\ref{o4r:verdict}, identifying the fibre counting
object as the Euler characteristic, with the corner candidate
refuted at $n = 3$ by the Morse lower bound alone. One finite check
remains from Section~\ref{sec:six}: the second-order block-form
verification of the warp-elimination identities recorded after
Proposition~\ref{o3:pairwise}.

Three hypotheses are structural. OW (Hypothesis~\ref{rb2:ow})
carries every lower bound of the topology; TR
(Hypothesis~\ref{rb2:tr}) carries the ring relation of the
projective generator; PS (Hypothesis~\ref{o2:ps}) carries the
Morse-theoretic admissibility, and its Palais--Smale clause is
expected to be the point of resistance. HS (Hypothesis~\ref{hyp:hs})
holds generically and suppresses only a failure of genericity at a
multiple coincidence. UB (Hypothesis~\ref{hyp:ub}) is the upper
bound on the critical count and is the hardest statement named in
this paper; the local route of Section~\ref{sec:retreat} exists
precisely so that the uniqueness consequence does not wait on it,
at the price of two local hypotheses B and U whose attack routes
are recorded where they are stated.

The division of labour between the two routes is worth one final
sentence. The structural route asserts that the topology of the
admissible region is the topology of the total space at generator
level, and derives the critical structure from it; the local route
asserts two rigidity statements at the constant stratum and derives
the uniqueness alone. Either route alone suffices for the statement
that the stable critical point is unique; only the structural route
delivers the count at every index; and the two consume disjoint
hypothesis sets, so a failure of either leaves the other standing.

\subsection*{Acknowledgements}

The author used Claude (Anthropic) and Gemini (Google DeepMind) as
mathematical development and analytical tools during the
preparation of this work. All intellectual content, theoretical
framework, and scientific claims are the sole responsibility of the
author.

\begin{thebibliography}{99}

\bibitem{vonNeumannWigner1929}
J. von Neumann and E. Wigner,
\emph{\"Uber das Verhalten von Eigenwerten bei adiabatischen
Prozessen},
Phys. Z. \textbf{30} (1929) 467--470.

\bibitem{Berry1984}
M. V. Berry,
\emph{Quantal phase factors accompanying adiabatic changes},
Proc. R. Soc. Lond. A \textbf{392} (1984) 45--57.

\bibitem{Simon1983}
B. Simon,
\emph{Holonomy, the quantum adiabatic theorem, and Berry's phase},
Phys. Rev. Lett. \textbf{51} (1983) 2167.

\bibitem{DysonMehta1963}
F. J. Dyson and M. L. Mehta,
\emph{Statistical theory of the energy levels of complex systems.
IV},
J. Math. Phys. \textbf{4} (1963) 701--712.

\bibitem{Kato}
T. Kato,
\emph{Perturbation Theory for Linear Operators},
Grundlehren der mathematischen Wissenschaften \textbf{132},
Springer, Berlin, 1995.

\bibitem{Hatcher2002}
A. Hatcher,
\emph{Algebraic Topology},
Cambridge University Press, Cambridge, 2002.

\bibitem{Palais1963}
R. S. Palais,
\emph{Morse theory on Hilbert manifolds},
Topology \textbf{2} (1963) 299--340.

\bibitem{Palais1979}
R. S. Palais,
\emph{The principle of symmetric criticality},
Comm. Math. Phys. \textbf{69} (1979) 19--30.

\bibitem{Milnor1963}
J. Milnor,
\emph{Morse Theory},
Annals of Mathematics Studies \textbf{51},
Princeton University Press, Princeton, 1963.

\bibitem{BottTu1982}
R. Bott and L. W. Tu,
\emph{Differential Forms in Algebraic Topology},
Graduate Texts in Mathematics \textbf{82},
Springer, New York, 1982.

\bibitem{Hitchin1974}
N. Hitchin,
\emph{Harmonic spinors},
Adv. Math. \textbf{14} (1974) 1--55.

\bibitem{Hijazi}
O. Hijazi,
\emph{A conformal lower bound for the smallest eigenvalue of the
Dirac operator and Killing spinors},
Comm. Math. Phys. \textbf{104} (1986) 151--162.

\bibitem{BourguignonGauduchon1992}
J.-P. Bourguignon and P. Gauduchon,
\emph{Spineurs, op\'erateurs de Dirac et variations de m\'etriques},
Comm. Math. Phys. \textbf{144} (1992) 581--599.

\bibitem{MEF3}
D. J. Master,
\emph{The Zwegers shadow as $U(1)$-projected intrinsic torsion on
twelve-manifolds},
Journal of Geometry and Physics, in press.
%% TODO-PI: complete bibliographic details (volume, year, DOI) of
%% the refereed Paper 3 before submission.

\end{thebibliography}

\end{document}
